\clearpage
\section{Altar Server Quick-Reference Handling Chart}
A category-ordered summary of all forty sanctuary objects, distinguishing what an altar server prepares, carries, or presents from what remains reserved to clergy or instituted ministers.

{\footnotesize\begin{longtable}{>{\raggedright\arraybackslash}p{0.23\linewidth}>{\raggedright\arraybackslash}p{0.27\linewidth}>{\raggedright\arraybackslash}p{0.43\linewidth}}
\toprule
\textbf{Object} & \textbf{Boundary} & \textbf{Practical rule} \\
\midrule
\endhead
\textbf{Altar} & Server recognizes only; does not handle & Servers venerate the altar with a profound bow on approaching or departing (genuflecting instead if the Blessed Sacrament is reserved in the sanctuary), place the Roman Missal and unconsecrated vessels upon it at the Preparation of the Gifts when needed, and never lean upon, clutter, or kiss the mensa (kissing the altar belongs to the bishop, priest, and deacon). \\
\textbf{Ambo} & Server recognizes only; does not handle & Servers flank the ambo with lighted candles and incense during the Gospel proclamation and may confirm before Mass that the Lectionary rests open upon the desk, but never proclaim readings or lead prayers from the ambo unless separately deputed as readers. \\
\textbf{Presidential Chair} & Clergy-reserved; server does not touch & Reserved strictly to the priest celebrant (or bishop when celebrating outside his cathedral); altar servers never sit in the presidential chair and approach it solely to hold the Roman Missal or Book of the Chair during prayers. \\
\textbf{Credence Table} & Server prepares or sets in place before or during Mass & Servers help arrange clean unconsecrated vessels, linens, cruets, and lavabo equipment on the credence table before Mass, carry needed items to the altar at the Preparation of the Gifts, and bring the water cruet for purification without touching unpurified sacred vessels. \\
\textbf{Tabernacle} & Clergy-reserved; server does not touch & Strictly reserved to clergy (and to instituted acolytes or duly deputed extraordinary ministers solely for Communion distribution or reposition); altar servers never open the tabernacle, handle its key, or touch reserved vessels, and genuflect before it outside Mass and at the beginning and end of Mass when it stands in the sanctuary (GIRM 274). \\
\textbf{Sanctuary Lamp} & Server recognizes only; does not handle & Servers look to the burning sanctuary lamp as the unmistakable sign that Christ is reserved in the tabernacle and never light, move, or extinguish it during Mass; any maintenance outside Mass follows sacristan direction with fire-safety care. \\
\textbf{Baptismal Font} & Server recognizes only; does not handle & Servers walk in procession to the font with cross, candles, ritual book, or drying towels during Baptisms and the Easter Vigil, never pouring sacramental water or disposing of blessed baptismal water improperly. \\
\textbf{Chalice} & Server prepares or sets in place before or during Mass & Servers may place the empty, unconsecrated chalice on the credence table before Mass and carry it to the altar at the Preparation of the Gifts if no deacon or instituted acolyte serves (GIRM 139). After consecration, servers never touch a chalice holding the Precious Blood, never consume from it, and never purify it (GIRM 279; RS 119; USCCB Norms 53). \\
\textbf{Paten} & Server prepares or sets in place before or during Mass & Servers may bring the unconsecrated paten and hosts to the altar at the Preparation of the Gifts (GIRM 139) and hold the separate Communion-plate level under communicants' chins during distribution (GIRM 118c; RS 93). Servers never touch consecrated hosts on the paten and never purify either the paten or the Communion-plate after Communion (GIRM 278--279; RS 119). \\
\textbf{Ciborium} & Server prepares or sets in place before or during Mass & Servers may set empty, unconsecrated ciboria on the credence table before Mass and bring empty ciboria to the altar at the Preparation of the Gifts or Fraction Rite when directed (USCCB Norms 37). Servers never open the tabernacle, never carry a ciborium containing consecrated hosts, and never purify a ciborium after Communion (GIRM 279; RS 119). \\
\textbf{Pyx} & Server recognizes only; does not handle & Servers may hand an empty, clean pyx to a priest or deacon in the sacristy when requested, but never handle, carry, open, or purify a pyx containing a consecrated Host. \\
\textbf{Communion Cup} & Server prepares or sets in place before or during Mass & Servers may set empty, clean Communion cups and folded purificators on the credence table before Mass and carry empty cups to the altar at the Preparation of the Gifts before consecration. Once the Eucharistic Prayer begins, servers never touch, carry, consume from, or purify Communion cups (GIRM 279; RS 119; USCCB Norms 53). \\
\textbf{Cruets} & Server carries and presents to priest or deacon & Cruets are core altar-server equipment: servers prepare them cleanly on the credence table, carry them to the altar with stoppers removed and handles turned toward the priest or deacon, return them promptly off the altar mensa (GIRM 306), and bring the water cruet back for vessel purification after Communion. \\
\textbf{Flagon} & Server handles directly & Servers may set the flagon at the gifts table before Mass, receive the unconsecrated flagon from the priest or deacon at the foot of the altar, and remove the empty flagon promptly from the altar to the credence table before the Eucharistic Prayer begins. \\
\textbf{Altar Cloth} & Server prepares or sets in place before or during Mass & Servers may assist the sacristan in laying or smoothing a clean altar cloth before Mass with clean hands, keeping wax, ink, and unauthorized items off the mensa and reporting any spill of the Precious Blood immediately to the priest or deacon. \\
\textbf{Corporal} & Server prepares or sets in place before or during Mass & Servers may set a clean, folded corporal on the credence table before Mass and bring it to the altar at the Preparation of the Gifts when no deacon or instituted acolyte serves; after consecration, servers never shake, refold carelessly, or wash a used corporal, and move a folded corporal flat and level only when directed after purification. \\
\textbf{Purificator} & Server prepares or sets in place before or during Mass & Servers may set clean, unused purificators on the credence table before Mass and carry them with unconsecrated chalices at the Preparation of the Gifts; once used with the Precious Blood, servers never wipe sacred vessels with them, shake them, or place them in ordinary laundry. \\
\textbf{Pall} & Server prepares or sets in place before or during Mass & Servers may set a clean pall atop the unconsecrated chalice on the credence table and bring it to the altar at the Preparation of the Gifts; once the Eucharistic Prayer begins, servers do not place or remove the pall over a chalice containing the Precious Blood. \\
\textbf{Lavabo Towel} & Server carries and presents to priest or deacon & Servers prepare the lavabo towel on the credence table before Mass, carry it with the water pitcher and basin to the side of the altar at the Preparation of the Gifts, present it cleanly unfolded for the priest to dry his hands, bow before and after, and return all three items promptly to the credence table. \\
\textbf{Roman Missal} & Server carries and presents to priest or deacon & Servers check ribbon markers before Mass as directed, hold the open Roman Missal steady at chest or reading height for the priest at the presidential chair when no stand or Book of the Chair is used, place it cleanly on the altar at the Preparation of the Gifts, and remove it after purification when instructed. \\
\textbf{Lectionary for Mass} & Server prepares or sets in place before or during Mass & Servers may confirm before Mass that the Lectionary lies open to the appointed day's readings on the ambo and may return it to the sacristy after Mass; servers never carry the Lectionary in the entrance procession and never proclaim Scripture from it unless separately deputed as readers. \\
\textbf{Book of the Gospels} & Server recognizes only; does not handle & Altar servers never carry the Book of the Gospels in the entrance procession (reserved strictly to the deacon, or to an instituted or deputed lector if no deacon is present) and never lift it from the altar for the Gospel procession; servers escort it with lighted candles and the thurible, flank the ambo during proclamation, and may reverently return the closed volume to the sacristy after Mass. \\
\textbf{Book of the Chair} & Server carries and presents to priest or deacon & Servers hold the Book of the Chair open flat and steady at the priest's reading angle during the Introductory Rites, Collect, Universal Prayer conclusion, Prayer after Communion, and Final Blessing, and never place it on the altar in place of the complete Roman Missal. \\
\textbf{Missal Stand} & Server handles directly & Servers keep the Missal stand off the altar mensa before the Preparation of the Gifts, carry it to the altar alongside the Roman Missal at the Preparation of the Gifts, set it neatly to the left of the corporal without overlapping the linen, and remove it cleanly after purification or Mass. \\
\textbf{Processional Cross} & Server carries in procession or movement & The cross-bearer carries the processional cross upright with the figure of Christ facing forward, omits bowing and genuflecting while holding it (GIRM 274), and sets it securely in its floor stand near the altar or puts it away if another altar cross already stands there (GIRM 122). \\
\textbf{Altar Candles} & Server handles directly & Servers light altar candles shortly before Mass, extinguish them cleanly with a snuffer afterward, and carry processional and Gospel candles upright at matching heights without bowing or genuflecting while holding them (GIRM 120b, 133, 274). \\
\textbf{Sanctuary Bells} & Server handles directly & An altar server rings the sanctuary bell cleanly and reverently shortly before the consecration and at each elevation according to local parish custom (GIRM 150), setting it down quietly without rattling the cups. \\
\textbf{Aspergillum and Aspersorium} & Server carries and presents to priest or deacon & Servers fill the clean vessel on the credence table before Mass, bring it to the priest for the blessing and sprinkling rite, accompany him through the nave holding the bucket steady, and return both articles cleanly to the credence table afterward (GIRM 51, 118c). \\
\textbf{Thurible} & Server handles directly & The thurifer lights charcoal safely in the sacristy, leads the entrance procession with the smoking thurible, raises the dome steadily for the priest to impose incense, hands the censer to the priest or deacon, and---when no deacon serves---incenses the priest, assembly, and consecrated Host and Chalice at the elevations (GIRM 120a, 150, 277). \\
\textbf{Incense Boat} & Server carries and presents to priest or deacon & The boat-bearer (or lone thurifer) opens the hinged lid and offers the spoon handle toward the priest or deacon before each imposition of incense, closing the lid promptly before walking so resin grains do not spill on sanctuary steps (GIRM 277). \\
\textbf{Alb} & Server handles directly & Altar servers may wear a clean white alb hemmed at ankle height (tied with a cincture if unfitted) where parish custom appoints it under US GIRM 339, hanging it neatly after Mass and never wearing clerical insignia over it (RS 126). \\
\textbf{Cincture} & Server handles directly & Altar servers tie the cincture neatly around the waist whenever wearing an unfitted alb, keeping the tasseled ends even and well clear of the floor and never confusing a colored waist cord with an ordained stole. \\
\textbf{Amice} & Server handles directly & Servers wearing a wide-necked alb that exposes ordinary shirt collars put on an amice beforehand (or use a high-collared alb that fully conceals street attire per GIRM 336), and assist in vesting a priest or bishop in the sacristy only when requested. \\
\textbf{Stole} & Clergy-reserved; server does not touch & Strictly reserved to bishops, priests, and deacons; altar servers, instituted acolytes, instituted lectors, and extraordinary ministers of Holy Communion never wear a stole (RS 126). A sacristan lays out the matching seasonal priestly or diaconal stole cleanly in the sacristy. \\
\textbf{Chasuble} & Clergy-reserved; server does not touch & Strictly reserved to bishops and priests celebrating or concelebrating Mass (GIRM 337); deacons, instituted acolytes, instituted lectors, extraordinary ministers, and altar servers never wear a chasuble (RS 126). During Mass, servers keep lighted candles and hot thuribles safely clear of its wide flowing drape. \\
\textbf{Dalmatic} & Clergy-reserved; server does not touch & Strictly reserved to deacons (and bishops beneath the chasuble at major pontifical celebrations); never worn by instituted acolytes, instituted lectors, extraordinary ministers, or altar servers (RS 126). Servers recognize the sleeved dalmatic immediately to distinguish a deacon from a priest. \\
\textbf{Cassock and Surplice} & Server handles directly & Altar servers wear cassock and surplice where parish custom appoints them (US GIRM 339), always wearing the surplice over a properly fitted cassock (never over street clothes alone) and never adding clerical collars, stoles, or prelatial insignia. \\
\textbf{Cope} & Clergy-reserved; server does not touch & Strictly reserved to bishops, priests, and deacons presiding at proper non-Mass rites; servers never wear a cope, though servers may assist in vesting it in the sacristy and---when directed during Benediction or processions---hold back its front edges as the celebrant incenses or blesses. \\
\textbf{Humeral Veil} & Server carries and presents to priest or deacon & Servers never wear the Eucharistic humeral veil to carry or bless with the Blessed Sacrament. During Benediction or Holy Thursday transfer, a server places the humeral veil evenly across the priest's or deacon's shoulders, fastens the front clasp securely, and removes and folds it cleanly afterward. \\
\textbf{Holy Oils and Oil Stocks} & Server moves only at celebrant's or sacristan's direction & Carry closed oil urns in procession only when designated (such as at the Chrism Mass or Holy Thursday reception of the oils) and hold a service tray with cotton and towel for the minister during Baptism or Confirmation; never open a stock, anoint anyone, transpose caps, or handle the ambry key. \\
\bottomrule
\end{longtable}}

\clearpage
\section{Pronunciation of Commonly Spoken Latin Terms}
Ecclesiastical (Roman) Latin pronunciation for Latin sanctuary headwords. Capitalized syllables take the primary stress.

{\footnotesize\begin{longtable}{>{\raggedright\arraybackslash}p{0.28\linewidth}>{\raggedright\arraybackslash}p{0.34\linewidth}>{\raggedright\arraybackslash}p{0.31\linewidth}}
\toprule
\textbf{Latin term} & \textbf{Ecclesiastical pronunciation} & \textbf{English headword} \\
\midrule
\endhead
\textit{Abacus} & AH-bah-koos & Credence Table \\
\textit{Alba} & AHL-bah & Alb \\
\textit{Altare} & ahl-TAH-reh & Altar \\
\textit{Ambo} & AHM-boh & Ambo \\
\textit{Amictus} & ah-MEEK-toos & Amice \\
\textit{Aspergillum et aspersorium} & ahs-pehr-JEEL-loom eht ahs-pehr-SOH-ree-oom & Aspergillum and Aspersorium \\
\textit{Calix} & KAH-leeks & Chalice \\
\textit{Campanula} & kahm-PAH-noo-lah & Sanctuary Bells \\
\textit{Candelabra cum cereis} & kahn-deh-LAH-brah koom CHEH-reh-ees & Altar Candles \\
\textit{Casula} & KAH-zoo-lah & Chasuble \\
\textit{Cingulum} & CHEEN-goo-loom & Cincture \\
\textit{Corporale} & kor-poh-RAH-leh & Corporal \\
\textit{Crux processionalis} & krooks proh-ches-see-oh-NAH-lees & Processional Cross \\
\textit{Dalmatica} & dahl-MAH-tee-kah & Dalmatic \\
\textit{Evangeliarium} & eh-vahn-jeh-lee-AH-ree-oom & Book of the Gospels \\
\textit{Fons baptismalis} & fohns bahp-tees-MAH-lees & Baptismal Font \\
\textit{Lagoena} & lah-JEH-nah & Flagon \\
\textit{Lampas peculiaris} & LAHM-pahs peh-koo-lee-AH-rees & Sanctuary Lamp \\
\textit{Lectionarium} & lehk-tsee-oh-NAH-ree-oom & Lectionary for Mass \\
\textit{Manutergium} & mah-noo-TEHR-jee-oom & Lavabo Towel \\
\textit{Mappa} & MAHP-pah & Altar Cloth \\
\textit{Missale Romanum} & mees-SAH-leh roh-MAH-noom & Roman Missal \\
\textit{Navicula} & nah-VEE-koo-lah & Incense Boat \\
\textit{Olea sancta et vasa oleorum} & OH-leh-ah SAHNK-tah eht VAH-sah oh-leh-OH-room & Holy Oils and Oil Stocks \\
\textit{Palla} & PAHL-lah & Pall \\
\textit{Patena} & pah-TEH-nah & Paten \\
\textit{Pluviale} & ploo-vee-AH-leh & Cope \\
\textit{Pulpitum} & POOL-pee-toom & Missal Stand \\
\textit{Purificatorium} & poo-ree-fee-kah-TOH-ree-oom & Purificator \\
\textit{Pyxis} & PEEK-sees & Ciborium \\
\textit{Pyxis parva} & PEEK-sees PAHR-vah & Pyx \\
\textit{Sedes sacerdotis celebrantis} & SEH-dehs sah-chehr-DOH-tees cheh-leh-BRAHN-tees & Presidential Chair \\
\textit{Stola} & STOH-lah & Stole \\
\textit{Tabernaculum} & tah-behr-NAH-koo-loom & Tabernacle \\
\textit{Turibulum} & too-REE-boo-loom & Thurible \\
\textit{Urceoli} & oor-CHEH-oh-lee & Cruets \\
\textit{Velum humerale} & VEH-loom oo-meh-RAH-leh & Humeral Veil \\
\textit{Vestis talaris et superpelliceum} & VEHS-tees tah-LAH-rees eht soo-pehr-pehl-LEE-cheh-oom & Cassock and Surplice \\
\bottomrule
\end{longtable}}

\clearpage
\section{English--Latin and Latin--English Indexes}
\subsection*{English to Latin}
{\footnotesize\begin{longtable}{>{\raggedright\arraybackslash}p{0.31\linewidth}>{\raggedright\arraybackslash}p{0.31\linewidth}>{\raggedright\arraybackslash}p{0.23\linewidth}>{\raggedright\arraybackslash}p{0.08\linewidth}}
\toprule
\textbf{English} & \textbf{Latin} & \textbf{Location} & \textbf{Note} \\
\midrule
\endhead
Alb & \textit{Alba} & Plate VIII & N31 \\
Altar & \textit{Altare} & Plate II & N01 \\
Altar Candles & \textit{Candelabra cum cereis} & Plate I & N26 \\
Altar Cloth & \textit{Mappa} & Plate IV & N15 \\
Ambo & \textit{Ambo} & Plate II & N02 \\
Amice & \textit{Amictus} & Plate VIII & N33 \\
Aspergillum and Aspersorium & \textit{Aspergillum et aspersorium} & Plate VI & N28 \\
Baptismal Font & \textit{Fons baptismalis} & Plate X & N07 \\
Book of the Chair & --- (no Latin typical edition) & Plate V & N23 \\
Book of the Gospels & \textit{Evangeliarium} & Plate V & N22 \\
Cassock and Surplice & \textit{Vestis talaris et superpelliceum} & Plate IX & N37 \\
Chalice & \textit{Calix} & Plate III & N08 \\
Chasuble & \textit{Casula} & Plate VIII & N35 \\
Ciborium & \textit{Pyxis} & Plate III & N10 \\
Cincture & \textit{Cingulum} & Plate VIII & N32 \\
Communion Cup & \textit{Calix} & Plate III & N12 \\
Cope & \textit{Pluviale} & Plate IX & N38 \\
Corporal & \textit{Corporale} & Plate IV & N16 \\
Credence Table & \textit{Abacus} & Plate I & N04 \\
Cruets & \textit{Urceoli} & Plate VI & N13 \\
Dalmatic & \textit{Dalmatica} & Plate VIII & N36 \\
Flagon & \textit{Lagoena} & Plate III & N14 \\
Holy Oils and Oil Stocks & \textit{Olea sancta et vasa oleorum} & Plate X & N40 \\
Humeral Veil & \textit{Velum humerale} & Plate IX & N39 \\
Incense Boat & \textit{Navicula} & Plate VII & N30 \\
Lavabo Towel & \textit{Manutergium} & Plate IV & N19 \\
Lectionary for Mass & \textit{Lectionarium} & Plate V & N21 \\
Missal Stand & \textit{Pulpitum} & Plate V & N24 \\
Pall & \textit{Palla} & Plate IV & N18 \\
Paten & \textit{Patena} & Plate III & N09 \\
Presidential Chair & \textit{Sedes sacerdotis celebrantis} & Plate II & N03 \\
Processional Cross & \textit{Crux processionalis} & Plate VI & N25 \\
Purificator & \textit{Purificatorium} & Plate IV & N17 \\
Pyx & \textit{Pyxis parva} & Plate III & N11 \\
Roman Missal & \textit{Missale Romanum} & Plate V & N20 \\
Sanctuary Bells & \textit{Campanula} & Plate VI & N27 \\
Sanctuary Lamp & \textit{Lampas peculiaris} & Plate II & N06 \\
Stole & \textit{Stola} & Plate VIII & N34 \\
Tabernacle & \textit{Tabernaculum} & Plate II & N05 \\
Thurible & \textit{Turibulum} & Plate VII & N29 \\
\bottomrule
\end{longtable}}

\subsection*{Latin to English}
{\footnotesize\begin{longtable}{>{\raggedright\arraybackslash}p{0.31\linewidth}>{\raggedright\arraybackslash}p{0.31\linewidth}>{\raggedright\arraybackslash}p{0.23\linewidth}>{\raggedright\arraybackslash}p{0.08\linewidth}}
\toprule
\textbf{Latin} & \textbf{English} & \textbf{Location} & \textbf{Note} \\
\midrule
\endhead
\textit{Abacus} & Credence Table & Plate I & N04 \\
\textit{Alba} & Alb & Plate VIII & N31 \\
\textit{Altare} & Altar & Plate II & N01 \\
\textit{Ambo} & Ambo & Plate II & N02 \\
\textit{Amictus} & Amice & Plate VIII & N33 \\
\textit{Aspergillum et aspersorium} & Aspergillum and Aspersorium & Plate VI & N28 \\
\textit{Calix} & Chalice & Plate III & N08 \\
\textit{Calix} & Communion Cup & Plate III & N12 \\
\textit{Campanula} & Sanctuary Bells & Plate VI & N27 \\
\textit{Candelabra cum cereis} & Altar Candles & Plate I & N26 \\
\textit{Casula} & Chasuble & Plate VIII & N35 \\
\textit{Cingulum} & Cincture & Plate VIII & N32 \\
\textit{Corporale} & Corporal & Plate IV & N16 \\
\textit{Crux processionalis} & Processional Cross & Plate VI & N25 \\
\textit{Dalmatica} & Dalmatic & Plate VIII & N36 \\
\textit{Evangeliarium} & Book of the Gospels & Plate V & N22 \\
\textit{Fons baptismalis} & Baptismal Font & Plate X & N07 \\
\textit{Lagoena} & Flagon & Plate III & N14 \\
\textit{Lampas peculiaris} & Sanctuary Lamp & Plate II & N06 \\
\textit{Lectionarium} & Lectionary for Mass & Plate V & N21 \\
\textit{Manutergium} & Lavabo Towel & Plate IV & N19 \\
\textit{Mappa} & Altar Cloth & Plate IV & N15 \\
\textit{Missale Romanum} & Roman Missal & Plate V & N20 \\
\textit{Navicula} & Incense Boat & Plate VII & N30 \\
\textit{Olea sancta et vasa oleorum} & Holy Oils and Oil Stocks & Plate X & N40 \\
\textit{Palla} & Pall & Plate IV & N18 \\
\textit{Patena} & Paten & Plate III & N09 \\
\textit{Pluviale} & Cope & Plate IX & N38 \\
\textit{Pulpitum} & Missal Stand & Plate V & N24 \\
\textit{Purificatorium} & Purificator & Plate IV & N17 \\
\textit{Pyxis} & Ciborium & Plate III & N10 \\
\textit{Pyxis parva} & Pyx & Plate III & N11 \\
\textit{Sedes sacerdotis celebrantis} & Presidential Chair & Plate II & N03 \\
\textit{Stola} & Stole & Plate VIII & N34 \\
\textit{Tabernaculum} & Tabernacle & Plate II & N05 \\
\textit{Turibulum} & Thurible & Plate VII & N29 \\
\textit{Urceoli} & Cruets & Plate VI & N13 \\
\textit{Velum humerale} & Humeral Veil & Plate IX & N39 \\
\textit{Vestis talaris et superpelliceum} & Cassock and Surplice & Plate IX & N37 \\
\bottomrule
\end{longtable}}

\clearpage
\section{Keyed Source Notes}
Each bracketed key on the plates resolves to the governing edition loci and 1962 comparison witnesses below.

{\footnotesize\begin{longtable}{>{\raggedright\arraybackslash}p{0.08\linewidth}>{\raggedright\arraybackslash}p{0.24\linewidth}>{\raggedright\arraybackslash}p{0.61\linewidth}}
\toprule
\textbf{Key} & \textbf{Entry} & \textbf{Governing loci and 1962 comparison} \\
\midrule
\endhead
\textbf{N01} & \textbf{Altar} & GIRM 273--274, 280, 296--308; USCCB GIRM 301, 304; CIC (1983) cann. 1235--1239; CDWDS Response Prot. N. 2036/00/L (25 Sept 2000). 1962 counterpart: Missale Romanum (1962), Rubricae Generales 526--527; Ritus Servandus II.2, IV.5, V.3. Symbolism: GIRM 296, 298; 1 Pet 2:4; 1 Cor 10:4; Heb 13:10; St. Thomas Aquinas, Summa Theologiae III, q. 83, a. 3. \\
\textbf{N02} & \textbf{Ambo} & GIRM 28, 118b, 130--136, 175, 196, 309; Ordo lectionum Missae (1981) Praenotanda 32--34; Inter Oecumenici (1964) 96. 1962 counterpart: Missale Romanum (1962), Ritus Servandus VI.4--5 (no fixed sanctuary ambo prescribed). Symbolism: GIRM 28, 309; Ordo lectionum Missae (1981) Praenotanda 32--34; Ps 118:105 [Heb 119:105]. \\
\textbf{N03} & \textbf{Presidential Chair} & GIRM 118a, 124, 127--130, 165--167, 310; Inter Oecumenici (1964) 92; Caeremoniale Episcoporum (1984). 1962 counterpart: Missale Romanum (1962), Ritus Servandus VI.4; Caeremoniale Episcoporum (1948) I.12--13; Fortescue (1943), pp. 8--9. Symbolism: GIRM 310. \\
\textbf{N04} & \textbf{Credence Table} & GIRM 118c, 139, 163, 178, 190, 192, 279, 306; Redemptionis Sacramentum (2004) 119; USCCB Norms for Communion under Both Kinds (2011) 53. 1962 counterpart: Missale Romanum (1962), Rubricae Generales 528; Ritus Servandus II.5; Caeremoniale Episcoporum (1948) I.12. \\
\textbf{N05} & \textbf{Tabernacle} & GIRM 274, 314--317; CIC (1983) can. 938; Redemptionis Sacramentum (2004) 130--131; Eucharisticum Mysterium (1967) 52--55. 1962 counterpart: Codex Iuris Canonici (1917) cann. 1268--1269; Missale Romanum (1962), Ritus Servandus II.2. Symbolism: GIRM 314--315; Eucharisticum Mysterium (1967) 53--55; Ex 25:8; St. Thomas Aquinas, Summa Theologiae III, q. 83, a. 5, ad 11. \\
\textbf{N06} & \textbf{Sanctuary Lamp} & GIRM 316; CIC (1983) can. 940. 1962 counterpart: Codex Iuris Canonici (1917) can. 1271. Symbolism: GIRM 316; CIC (1983) can. 940; Ex 27:20--21; Lev 24:2--3. \\
\textbf{N07} & \textbf{Baptismal Font} & CIC (1983) can. 858, can. 861 §1; Ordo Baptismi Parvulorum (General Introduction 25, 45--47); Roman Missal (2008/2011) Easter Vigil Blessing of Baptismal Water; CCC 1237--1243. 1962 counterpart: Missale Romanum (1962), Vigilia Paschalis (cf. Missale Romanum 1922, Holy Saturday); Rituale Romanum (1925/1952), Tit. II. Symbolism: Roman Missal Easter Vigil; CCC 1237--1243; John 3:5; Rom 6:3--4; St. Cyril of Jerusalem, Mystagogical Catechesis 2.4. \\
\textbf{N08} & \textbf{Chalice} & GIRM 118c, 139, 142, 150--151, 163, 279, 306, 327--330, 333; Redemptionis Sacramentum 106, 117--119; USCCB Norms for Holy Communion under Both Kinds (2011) 35, 52--55. 1962 counterpart: Missale Romanum (1962), Ritus Servandus I.1, II.1, VIII. Symbolism: GIRM 327; Matt 26:27--28; 1 Cor 10:16; Ps 115:13 [Heb 116:13]; St. Thomas Aquinas, Summa Theologiae III, q. 83, a. 3, ad 6 \& a. 5, ad 9. \\
\textbf{N09} & \textbf{Paten} & GIRM 118c, 139, 155, 278--279, 306, 327--329, 331, 333; USCCB GIRM 329; Redemptionis Sacramentum 93, 117--119. 1962 counterpart: Missale Romanum (1962), Ritus Servandus I.1, VII.9, X.1--2, X.7. Symbolism: GIRM 327, 331; 1 Cor 10:17; St. Thomas Aquinas, Summa Theologiae III, q. 83, a. 3, obj. 6 \& ad 6. \\
\textbf{N10} & \textbf{Ciborium} & GIRM 118c, 163, 279, 306, 327--329, 333; USCCB GIRM 329; CIC (1983) can. 939; Redemptionis Sacramentum 117--119; USCCB Norms for Holy Communion under Both Kinds (2011) 37, 53. 1962 counterpart: Missale Romanum (1962), Ritus Servandus II.3, VII.3, X.6; CIC (1917) can. 1269. \\
\textbf{N11} & \textbf{Pyx} & GIRM 327--329, 333; USCCB GIRM 329; CIC (1983) can. 939; Redemptionis Sacramentum 117--119, 133; Holy Communion and Worship of the Eucharist Outside Mass (1973). 1962 counterpart: Rituale Romanum (1925/1952), Title IV, ch. 2; Missale Romanum (1962), Ritus Servandus X.6. \\
\textbf{N12} & \textbf{Communion Cup} & GIRM 279, 281--287, 327--330, 333; USCCB GIRM 283; USCCB Norms for Holy Communion under Both Kinds (2011) 32--38, 52--55; Redemptionis Sacramentum 100--107, 117--119. 1962 counterpart: Missale Romanum (1962), Rubricae Generales 503; Ritus Servandus X.6. Symbolism: GIRM 281; USCCB Norms for Holy Communion under Both Kinds (2011) 32--33; Matt 26:27--28; 1 Cor 10:16; 1 Cor 11:25--26. \\
\textbf{N13} & \textbf{Cruets} & GIRM 118c, 139--142, 163, 178, 190, 279, 306; CIC (1983) can. 924 §3. 1962 counterpart: Missale Romanum (1962), Rubricae Generales 528; Ritus Servandus VII.4, VII.6, VII.10. \\
\textbf{N14} & \textbf{Flagon} & GIRM 73, 118c, 140, 306, 328--330; Redemptionis Sacramentum 106, 117; USCCB Norms for Holy Communion under Both Kinds (2011) 35--38. 1962 counterpart: Missale Romanum (1962), Rubricae Generales 528 (negative witness: small wine and water cruets only). \\
\textbf{N15} & \textbf{Altar Cloth} & GIRM 117, 304 (with USCCB territorial adaptation); GIRM 280, 306; Redemptionis Sacramentum 120. 1962 counterpart: Missale Romanum (1962), Rubricae Generales 526. Symbolism: GIRM 304; St. Thomas Aquinas, Summa Theologiae III, q. 83, a. 3, ad 7; Matt 27:59. \\
\textbf{N16} & \textbf{Corporal} & GIRM 118c, 139, 163, 178, 183, 190, 192, 279--280, 306; Redemptionis Sacramentum 120; USCCB Norms for Holy Communion under Both Kinds (2011) 53. 1962 counterpart: Missale Romanum (1962), Ritus Servandus I.1, II.2, VI.7. Symbolism: St. Thomas Aquinas, Summa Theologiae III, q. 83, a. 3, ad 7 (citing De Consecr. dist. 1); Matt 27:59; John 20:6. \\
\textbf{N17} & \textbf{Purificator} & GIRM 118c, 139, 163, 178, 192, 279, 286, 306; Redemptionis Sacramentum 119--120; USCCB Norms for Holy Communion under Both Kinds (2011) 43, 53. 1962 counterpart: Missale Romanum (1962), Ritus Servandus I.1, VII.4, X.5. \\
\textbf{N18} & \textbf{Pall} & GIRM 118c (pro opportunitate), 139, 306; Redemptionis Sacramentum 120. 1962 counterpart: Missale Romanum (1962), Ritus Servandus I.1, VII.5, VIII.7. \\
\textbf{N19} & \textbf{Lavabo Towel} & GIRM 76, 118c, 145, 190. 1962 counterpart: Missale Romanum (1962), Rubricae Generales 528; Ritus Servandus VII.6--7, VII.10; Ordo Missae 1039. \\
\textbf{N20} & \textbf{Roman Missal} & GIRM 30, 118a, 139, 147, 189, 306, 349, 389; Postquam Summus Pontifex (2021) 2--3, 40. 1962 counterpart: Missale Romanum (1962), Rubricae Generales 527; Ritus Servandus I.1, II.4--5, VI.5, XI.1. Symbolism: GIRM 349. \\
\textbf{N21} & \textbf{Lectionary for Mass} & GIRM 28, 118b, 120d, 172, 194, 349; Ordo lectionum Missae (1981) Praenotanda 35--37, 69, 103--107. 1962 counterpart: Missale Romanum (1962), Ritus Servandus VI.4. Symbolism: GIRM 28, 349; Ordo lectionum Missae (1981) Praenotanda 35; Ps 118:105 [119:105]. \\
\textbf{N22} & \textbf{Book of the Gospels} & GIRM 62, 117, 120d, 122, 133--134, 172--175, 194--195, 277, 306, 349; Ordo lectionum Missae (1981) Praenotanda 36--37. 1962 counterpart: Missale Romanum (1962), Ritus Servandus VI.5. Symbolism: GIRM 62, 134, 175, 349; St. Thomas Aquinas, Summa Theologiae III, q. 83, a. 4, corpus. \\
\textbf{N23} & \textbf{Book of the Chair} & GIRM 118a, 306; USCCB Excerpts from the Roman Missal (CDWDS confirmed, 2018). 1962 counterpart: No presbyteral counterpart in Missale Romanum (1962), Rubricae Generales 527; episcopal book-bearer precedent only in Caeremoniale Episcoporum (1948) I.20. \\
\textbf{N24} & \textbf{Missal Stand} & GIRM 139, 306--307 (practical custom supporting the altar Missal). 1962 counterpart: Missale Romanum (1962), Rubricae Generales 527; Ritus Servandus VI.5, XI.1. \\
\textbf{N25} & \textbf{Processional Cross} & GIRM 117; GIRM 120b; GIRM 122; GIRM 188; GIRM 274; GIRM 308. 1962 counterpart: Missale Romanum (1962), Rubricae Generales n. 527; Fortescue (1943), pp. 17, 71, 78--79. Symbolism: GIRM 308; 1 Cor 1:23; Gal 6:14; Col 1:20; St. Thomas Aquinas, Summa Theologiae III, q. 83, a. 5, ad 3. \\
\textbf{N26} & \textbf{Altar Candles} & GIRM 117; GIRM 120b; GIRM 122; GIRM 133; GIRM 175; GIRM 274; GIRM 307. 1962 counterpart: Missale Romanum (1962), Rubricae Generales n. 527; Ritus Servandus II.5, VI.5; Fortescue (1943), pp. 17, 89--90. Symbolism: GIRM 117, 307; John 8:12. \\
\textbf{N27} & \textbf{Sanctuary Bells} & GIRM 150. 1962 counterpart: Missale Romanum (1962), Rubricae Generales n. 528; Ritus Servandus VII.8, VIII.6, X.6. \\
\textbf{N28} & \textbf{Aspergillum and Aspersorium} & GIRM 51; GIRM 118c; Roman Missal (2008/2011) Appendix II; Tres abhinc annos (1967) 26. 1962 counterpart: Missale Romanum (1962), Ordo ad aspergendam aquam benedictam, pp. [231]--[233]; Fortescue (1917), p. 16; Fortescue (1943), pp. 79--81. Symbolism: GIRM 51; Roman Missal, Appendix II; Ps 50:9 [Heb 51:7]. \\
\textbf{N29} & \textbf{Thurible} & GIRM 120a; GIRM 133; GIRM 144; GIRM 150; GIRM 175; GIRM 179; GIRM 276; GIRM 277. 1962 counterpart: Missale Romanum (1962), Ritus Servandus IV.4, VI.5, VII.10, VIII.8; Ordo Missae nn. 1035--1038. Symbolism: GIRM 276; Ps 140:2 [Heb 141:2]; Apoc 5:8, 8:3--4; 2 Cor 2:14; St. Thomas Aquinas, Summa Theologiae III, q. 83, a. 5, ad 2. \\
\textbf{N30} & \textbf{Incense Boat} & GIRM 120; GIRM 144; GIRM 276; GIRM 277; Caeremoniale Episcoporum (1984). 1962 counterpart: Missale Romanum (1962), Ritus Servandus IV.4, VII.10. \\
\textbf{N31} & \textbf{Alb} & GIRM 119; GIRM 335; GIRM 336; GIRM 339 (USCCB); Redemptionis Sacramentum 121--122, 126; Caeremoniale Episcoporum (1984) 65. 1962 counterpart: Missale Romanum (1962), Rubricae Generales Tit. XIX.1; Ritus Servandus I.2--3, II.1, VI.8; Praeparatio ad Missam (Ad Albam). Symbolism: GIRM 335--336; Gal 3:27; Apoc 7:13--14; traditional Roman sacristy vesting prayer (Ad Albam); Amalarius of Metz, Liber Officialis II.18. \\
\textbf{N32} & \textbf{Cincture} & GIRM 336; GIRM 340; Redemptionis Sacramentum 121--122; Caeremoniale Episcoporum (1984) 65. 1962 counterpart: Missale Romanum (1962), Ritus Servandus I.3; Praeparatio ad Missam (Ad Cingulum). Symbolism: Practical norm: GIRM 336; devotional tradition: Roman sacristy vesting prayer (Ad Cingulum); Luke 12:35; Eph 6:14; Amalarius of Metz, Liber Officialis II.22. \\
\textbf{N33} & \textbf{Amice} & GIRM 336; Redemptionis Sacramentum 121--122; Caeremoniale Episcoporum (1984) 65. 1962 counterpart: Missale Romanum (1962), Ritus Servandus I.3; Praeparatio ad Missam (Ad Amictum). Symbolism: Practical norm: GIRM 336; devotional tradition: Roman sacristy vesting prayer (Ad Amictum); Eph 6:17; Ps 38:2 [Heb 39:1]; Ps 140:3 [Heb 141:3]; Amalarius of Metz, Liber Officialis II.17. \\
\textbf{N34} & \textbf{Stole} & GIRM 119; GIRM 336; GIRM 337; GIRM 338; GIRM 340; GIRM 346 (USCCB 346e, 346h); Redemptionis Sacramentum 123--126. 1962 counterpart: Missale Romanum (1962), Rubricae Generales n. 137; Ritus Servandus I.3--4; Praeparatio ad Missam (Ad Stolam). Symbolism: GIRM 335, 340; traditional Roman sacristy vesting prayer (Ad Stolam); Matt 11:29--30; Amalarius of Metz, Liber Officialis II.20. \\
\textbf{N35} & \textbf{Chasuble} & GIRM 119a; GIRM 209; GIRM 335; GIRM 337; GIRM 344; GIRM 346 (USCCB 346e, 346h); Redemptionis Sacramentum 123--124, 126--127; Tres abhinc annos (1967) 26. 1962 counterpart: Missale Romanum (1962), Rubricae Generales nn. 133--134, 137; Ritus Servandus I.4; Praeparatio ad Missam (Ad Casulam). Symbolism: GIRM 335, 337; Matt 11:30; Col 3:14; traditional Roman sacristy vesting prayer (Ad Casulam); Amalarius of Metz, Liber Officialis II.19. \\
\textbf{N36} & \textbf{Dalmatic} & GIRM 119b, 335--336, 338, 340, 342--346; Redemptionis Sacramentum 125--126; Caeremoniale Episcoporum (1984) 56. 1962 counterpart: Missale Romanum (1962), Rubricae Generales 134, 137; Caeremoniale Episcoporum (1948) II.8 n. 18. Symbolism: GIRM 335, 338; Caeremoniale Episcoporum (1984) 56; Is 61:10; James 1:27; traditional Roman vesting prayer; Amalarius of Metz, Liber Officialis II.21. \\
\textbf{N37} & \textbf{Cassock and Surplice} & GIRM 119c, 336, 339; USCCB GIRM 339; Redemptionis Sacramentum 128; Caeremoniale Episcoporum (1984) 65. 1962 counterpart: Missale Romanum (1962), Ritus Servandus I.2, II.1, VI.8; Caeremoniale Episcoporum (1948) I.11 §5. \\
\textbf{N38} & \textbf{Cope} & GIRM 337, 341, 346; Holy Communion and Worship of the Eucharist Outside Mass (1973) 91--92; Tres abhinc annos (1967) 26; Caeremoniale Episcoporum (1984). 1962 counterpart: Missale Romanum (1962), Rubricae Generales Cap. XIX (137); Caeremoniale Episcoporum (1948) II.8 n. 18. \\
\textbf{N39} & \textbf{Humeral Veil} & Holy Communion and Worship of the Eucharist Outside Mass (1973) 91--92; Roman Missal (2002/2008/2011), Holy Thursday Evening Mass of the Lord's Supper; Caeremoniale Episcoporum (1984); Ministeria quaedam (1972). 1962 counterpart: Missale Romanum (1962), Ritus Servandus VII.9, X.8; Holy Thursday Transfer. Symbolism: Holy Communion and Worship of the Eucharist Outside Mass (1973) nn. 91--92; Roman Missal (2008/2011), Evening Mass of the Lord's Supper. \\
\textbf{N40} & \textbf{Holy Oils and Oil Stocks} & CIC 1983 can. 847, pars. 1--2; CIC 1983 can. 999 n. 2; can. 1003, par. 1; Ordo benedicendi oleum catechumenorum et infirmorum et conficiendi chrisma (3 Dec 1970; editio typica 1971); Caeremoniale Episcoporum (14 Sept 1984) no. 274; Missale Romanum (2002/2008; US 2011), Missa Chrismatis; CCC nos. 1241, 1294--1300. 1962 counterpart: Missale Romanum (1962), Feria V in Coena Domini (Missa Chrismatis) and Vigilia Paschalis (Blessing of Baptismal Water); Pontificale Romanum (1859), De benedictione oleorum. Symbolism: CCC 1294--1300; Ordo benedicendi oleum (1970); Roman Chrism Preface; 2 Cor 1:21--22; James 5:14--15. \\
\bottomrule
\end{longtable}}

\clearpage
\section{Scope, Edition, and Qualifications}
\textbf{Governing edition and jurisdiction.} Entries in the altar server visual dictionary describe the Roman Rite as governed by the \textit{Missale Romanum}, \textit{editio typica tertia} (2002; \textit{reimpressio emendata}, 2008), and its \textit{Institutio Generalis Missalis Romani} (GIRM), together with the approved English edition for the United States of America, \textit{The Roman Missal, Third Edition} (effective First Sunday of Advent, 27 November 2011), the adaptations approved by the United States Conference of Catholic Bishops (USCCB), the 2021 emendations incorporated in the official USCCB GIRM (Prot.~n.~228/20; Prot.~n.~394/21, 22 October 2021), and the 2021 amendment of canon 230 \S1 (\textit{Spiritus Domini}, 10 January 2021). Norms are stated as of 2026-11-25.

\textbf{Coverage and boundaries.} Within this practical reference, 40 principal sanctuary furnishings, sacred vessels, linens, liturgical books, service items, and vestments of Mass and closely joined rites are documented for active parish service; historically superseded postconciliar forms and full 1962 Roman Missal comparisons are collected separately in the comprehensive adult reference edition. Episcopal and pontifical insignia (miter, crosier, pectoral cross, ring, zucchetto, gremiale, bugia), full Paschal Triduum apparatus (Paschal candle, Easter fire brazier, Good Friday adoration cross veil), funeral pall and catafalque, and specialized rites of dedication and ordination remain outside the present scope and are reserved for future expansion rather than treated by conjecture.

\textbf{Separation of universal law, territorial adaptation, and custom.} Universal provisions of the Latin typical edition (including choices entrusted to the Diocesan Bishop, such as tabernacle placement at GIRM 315), approved United States adaptations (solid wood fixed altars at GIRM 301, festive under-cloths at GIRM 304, sacred-vessel materials at GIRM 329, lay ministerial attire at GIRM 339, Communion under both kinds at GIRM 283 and the USCCB Norms, and funeral and solemn vestment colors at GIRM 346e and 346h), diocesan or celebrant options, and documented parish custom are distinguished throughout. Counterpart references to the 1962 Roman Missal are restricted to the Keyed Source Notes (with full side-by-side comparisons in the comprehensive adult reference edition) and are never imported as postconciliar norms.

\textbf{Rights, visual provenance, and reliance boundary.} All ten monochrome graphite plates are original project-created illustrations normalized to 8-bit grayscale on pure white ground; no third-party scans, photographs, or copyrighted illustrations are reproduced. English descriptions paraphrase official liturgical law in original prose without reproducing protected ICEL liturgical prayer texts or rubrical formulations verbatim. English Scripture quotations follow the public-domain Douay--Rheims (Challoner) text. The approved Roman Missal and diocesan directives remain the governing liturgical law; study and preparation notes here carry no official, liturgical, or magisterial authority.

\phantomsection
\section*{References}
\addcontentsline{toc}{section}{References}
\subsection*{Governing Postconciliar Roman Liturgical Books and Law}
\begin{itemize}
\item Missale Romanum, ex decreto Sacrosancti Oecumenici Concilii Vaticani II instauratum, auctoritate Pauli PP. VI promulgatum, Ioannis Pauli PP. II cura recognitum, editio typica tertia (Typis Vaticanis, 2002); reimpressio emendata (Typis Vaticanis / Libreria Editrice Vaticana, 2008; announced in Notitiae 44 [2008], nos. 503--504, pp. 367--388), including Institutio Generalis Missalis Romani (IGMR).
\item The Roman Missal, Third Edition, for Use in the Dioceses of the United States of America (Washington, DC: United States Conference of Catholic Bishops, 2011; CDWDS recognitio 26 March 2010 and 24 July 2010; effective 27 November 2011), including the General Instruction of the Roman Missal with U.S. adaptations and 2021 emendations (Prot. n. 228/20; Prot. n. 394/21, 22 October 2021).
\item United States Conference of Catholic Bishops, Norms for the Distribution and Reception of Holy Communion under Both Kinds in the Dioceses of the United States of America (approved 15 June 2001; CDWDS recognitio 22 March 2002, Prot. 1383/01/L; revised 2011; reflecting October 2006 CDWDS decision on vessel purification).
\item United States Conference of Catholic Bishops, Excerpts from the Roman Missal (Book of the Chair; CDWDS confirmed; Washington, DC: USCCB, 2018).
\item Congregation for Divine Worship and the Discipline of the Sacraments, Response on GIRM no. 299 (25 September 2000, Prot. N. 2036/00/L; Notitiae 36 [2000], 171--173); Instruction Redemptionis Sacramentum (25 March 2004; AAS 96 [2004], 549--601); Decree implementing the provisions of canon 838 of the Code of Canon Law, Postquam Summus Pontifex (22 October 2021).
\item Ordo lectionum Missae, editio typica altera (Libreria Editrice Vaticana, 21 January 1981), Praenotanda nos. 1--125.
\item Caeremoniale Episcoporum, ex decreto Sacrosancti Oecumenici Concilii Vaticani II instauratum, editio typica (Typis Polyglottis Vaticanis, 14 September 1984).
\item Codex Iuris Canonici auctoritate Ioannis Pauli PP. II promulgatus (25 January 1983; AAS 75/II [1983]), cann. 230 (as amended by Pope Francis, Motu Proprio Spiritus Domini, 10 January 2021), 847, 858, 910, 938--940, 1003, 1235--1239.
\item Catechism of the Catholic Church, 2nd ed. (Vatican City: Libreria Editrice Vaticana, 1997; Washington, DC: USCCB, 2019).
\item Sacred Congregation for Divine Worship, Ordo Baptismi Parvulorum (Rite of Baptism for Children; editio typica, 1969; editio typica altera, 1973); Ordo benedicendi oleum catechumenorum et infirmorum et conficiendi chrisma (3 December 1970); De sacra communione et de cultu mysterii eucharistici extra Missam (21 June 1973).
\end{itemize}
\subsection*{Earlier Postconciliar Legislative Milestones (1963--1975)}
\begin{itemize}
\item Second Vatican Council, Constitution on the Sacred Liturgy Sacrosanctum Concilium (4 December 1963; AAS 56 [1964], 97--138).
\item Sacred Congregation of Rites and Consilium, Instruction Inter Oecumenici (26 September 1964; AAS 56 [1964], 877--900), Instruction Tres abhinc annos (4 May 1967; AAS 59 [1967], 442--448), and Instruction Eucharisticum Mysterium (25 May 1967; AAS 59 [1967], 539--573).
\item Pope Paul VI, Apostolic Constitution Missale Romanum (3 April 1969; AAS 61 [1969], 217--222); Motu Proprio Mysterii Paschalis (14 February 1969; AAS 61 [1969], 222--226); Motu Proprio Ministeria quaedam (15 August 1972; AAS 64 [1972], 529--534).
\item Sacred Congregation for Divine Worship, Missale Romanum, editio typica prima (26 March 1970, Prot. 166/70) and editio typica altera (27 March 1975); Sacred Congregation for the Discipline of the Sacraments, Instruction Immensae caritatis (29 January 1973; AAS 65 [1973], 264--271).
\end{itemize}
\subsection*{1962 Comparison Witnesses, Scripture, and Liturgical Tradition}
\begin{itemize}
\item Missale Romanum, editio typica (Typis Polyglottis Vaticanis, 1962), including Rubricae Generales Missalis Romani (1960/1962) and Ritus Servandus in Celebratione Missae; together with Codex Iuris Canonici (1917), Rituale Romanum (1925/1952), and Pontificale Romanum (1859).
\item Adrian Fortescue, The Ceremonies of the Roman Rite Described (London: Burns \& Oates, 1917; 7th rev. ed., 1943).
\item The Holy Bible, Douay--Rheims Version (Challoner revision).
\item St. Cyril of Jerusalem, Mystagogical Catecheses (II.4); St. Thomas Aquinas, Summa Theologiae, Tertia Pars, q. 83, aa. 3--5 (on sacred vessels, linens, vestments, and altar rites); Amalarius of Metz, Liber Officialis (De ecclesiasticis officiis, PL 105); traditional Roman vesting prayers of the Missale Romanum.
\end{itemize}
